\section{Read bytes in the order the hash specifies}
The simulator's memory view and the hash's word interpretation answer different
questions. In the chapter's little-endian memory configuration, a word load from bytes
\texttt{61 62 63 80} yields the low word \texttt{80636261}. SHA-256 interprets those
same four bytes as the big-endian word \texttt{61626380}. We must construct the latter
value explicitly.

Load each byte with \texttt{lbu}, shift the first three left by 24, 16 and 8 bits,
and OR them with the fourth byte. The positions refer to the numeric word, not to a
change in the bytes held in the source. Byte loads also allow the block to begin at
an unaligned address. The destination schedule is a word array, so its address must
be word-aligned.

The first sixteen schedule words are just the sixteen groups of four bytes in the
block. For padded \texttt{abc}, \(W_0=\mathtt{61626380}\), words 1 through 14 are
zero, and \(W_{15}=\mathtt{00000018}\). After constructing a numeric word, store it
with \texttt{sw}. Reading that schedule word with \texttt{lw} recovers the low-word
pattern we constructed, even though its bytes are stored in little-endian order.
The source block and the schedule therefore have different storage roles.

A signed byte load can be particularly misleading in the final position. Suppose the
last byte is \texttt{80}. A signed load fills all higher bits with ones. ORing it
with the other three fields then changes the entire upper part of the word. Unsigned
loads keep each byte confined to its intended eight bits before combination.

\MARSCodeLabel{Decode one source word; advance both pointers after the store.}
\begin{minted}{mars-idea}
    lbu $t2, 0($a0)
    lbu $t3, 1($a0)
    sll $t2, $t2, 24
    sll $t3, $t3, 16
    or $t2, $t2, $t3
    lbu $t3, 2($a0)
    sll $t3, $t3, 8
    or $t2, $t2, $t3
    lbu $t3, 3($a0)
    or $t2, $t2, $t3
    sw $t2, 0($t0)
\end{minted}
Here \texttt{a0} points to the source group and \texttt{t0} points to its schedule
word. Advance both by four bytes after the store. Stop after sixteen words; the next
stage writes the remaining schedule entries from these decoded values.
